==============================================================================
            GUÍA DE ARQUITECTURA Y USO DEL PROYECTO LATEX (TFG)
==============================================================================

1. ESTRUCTURA DE CARPETAS
------------------------------------------------------------------------------
El proyecto sigue una arquitectura modular "Árbol". El archivo raíz es el director
y el contenido está ordenado en la carpeta 'recursos'.

PROYECTO/
├── main.tex                <-- NO TOCAR MUCHO. Solo orquesta los archivos.
├── LEEME.txt               <-- Este archivo.
└── recursos/               <-- AQUÍ ESTÁ TODO EL CONTENIDO REAL
    ├── capitulos/          <-- Tus capítulos (Introducción, Desarrollo...)
    ├── config/             <-- Configuración técnica
    │   ├── UAL.sty         <-- Diseño visual (Colores, Cabeceras, Portada)
    │   ├── datos.tex       <-- Título, Autores, Tutores, Curso
    │   ├── paquetes.tex    <-- Lista de \usepackage (Librerías)
    │   └── codigo.sty      <-- Configuración para mostrar código (Python)
    ├── frontmatter/        <-- Resumen e Índices
    ├── backmatter/         <-- Apéndices y Contraportada
    ├── Figuras/            <-- Todas tus imágenes van aquí
    └── referencias.bib     <-- Archivo de bibliografía

------------------------------------------------------------------------------
2. GUÍABURROS: ¿CÓMO HAGO... ?
------------------------------------------------------------------------------

[A] ... CAMBIAR MIS DATOS (AUTOR, TÍTULO)
    Ve a: recursos/config/datos.tex
    Ahí puedes cambiar los nombres de los alumnos, directores y título del TFG.
    Se actualizarán automáticamente en la portada.

[B] ... ESCRIBIR O AÑADIR UN CAPÍTULO NUEVO
    1. Crea un archivo .tex nuevo en: recursos/capitulos/ (ej: 04-resultados.tex).
    2. Escribe tu contenido dentro.
    3. Ve a 'main.tex' y busca la zona "CUERPO DEL TRABAJO".
    4. Añade la línea: \input{recursos/capitulos/04-resultados}

[C] ... CAMBIAR LOS COLORES DEL DOCUMENTO
    Ve a: recursos/config/UAL.sty
    Busca la sección "PALETA DE COLORES" para ver los disponibles.
    Baja a la sección "DISEÑO DE TÍTULOS" o "CABECERAS".
    Cambia el color (ej: 'BlueDark') por otro (ej: 'RedAccent').
    ¡Todo el documento cambiará automáticamente!

[D] ... INSERTAR UNA IMAGEN
    Sube la imagen a la carpeta: recursos/Figuras/
    En tu texto usa: \includegraphics[width=...]{nombre_imagen.png}
    IMPORTANTE: NO escribas la ruta "recursos/Figuras/...", solo el nombre.
    LaTeX ya sabe dónde buscar gracias al comando \graphicspath en el main.

[E] ... AÑADIR REFERENCIAS BIBLIOGRÁFICAS
    1. Busca la cita en Google Scholar -> Citar -> BibTeX.
    2. Pega el código en: recursos/referencias.bib
    3. En tu texto cita usando: \cite{clave_de_la_cita}

[F] ... PEGAR CÓDIGO (PYTHON/C)
    Usa el entorno lstlisting:
    \begin{lstlisting}[language=Python, caption={Descripción}]
       print("Hola Mundo")
    \end{lstlisting}

------------------------------------------------------------------------------
3. SOLUCIÓN DE PROBLEMAS COMUNES
------------------------------------------------------------------------------

* ERROR: "File `nombre_imagen.png` not found"
  -> Verifica mayúsculas/minúsculas. "Foto.png" no es igual a "foto.png".
  -> Verifica que la imagen esté dentro de recursos/Figuras.

* ERROR: La compilación se rompe tras borrar archivos temporales
  -> Overleaf a veces guarda "basura". Dale a "Recompile from scratch" (en el
     menú del botón verde de recompilar).

* ERROR: Los colores salen negros
  -> Verifica en UAL.sty que no has borrado la definición de colores al principio.

==============================================================================